\section{Random variables as functions on the sample space}

Until now, probability models have been built around sample spaces and events.
We described an experiment, decided what counts as an elementary outcome,
constructed a sample space \(\Omega\), and assigned probabilities to events. A
random variable adds a new layer: it extracts a numerical quantity from each
outcome.

This is the central idea:
\[
\text{a random variable is a function on the sample space.}
\]
It is not a mysterious number that is random by itself. Its randomness comes from
the randomness of the elementary outcome to which the function is applied.

\begin{definitionbox}
Let \(\Omega\) be a sample space. A \textbf{random variable} is a function
\[
X:\Omega\to\mathbb{R}.
\]
For each elementary outcome \(\omega\in\Omega\), the value \(X(\omega)\) is a
real number extracted from that outcome.
\end{definitionbox}

The notation may look simple, but it changes the way we organize probability.
Instead of asking only which event occurred, we can now ask which number was
produced by a function of the outcome.

\subsection*{Why introduce random variables?}

Many experiments produce outcomes that are richer than the numerical quantity we
want to study. For example, if a coin is tossed five times, one elementary
outcome may be
\[
\omega=HHTTH.
\]
The full outcome records the order of all tosses. But in many questions we care
only about the number of heads. Then we define a function
\[
X(\omega)=\text{the number of heads in }\omega.
\]
For this outcome,
\[
X(HHTTH)=3.
\]

The random variable does not replace the sample space. It is built on top of it.
The sample space contains detailed outcomes; the random variable extracts the
numerical feature we want to analyze.

\subsection*{Example: the sum of two dice}

Roll two fair dice and record the ordered pair. The sample space is
\[
\Omega=\{(i,j):i,j\in\{1,2,3,4,5,6\}\}.
\]
An elementary outcome is a pair \((i,j)\). If we are interested in the sum, we
define
\[
S:\Omega\to\mathbb{R},
\qquad
S(i,j)=i+j.
\]
Thus
\[
S(1,6)=7,\qquad S(3,4)=7,\qquad S(6,6)=12.
\]
Different elementary outcomes may give the same value of the random variable.
This is important. The value \(7\) is not one elementary outcome in the two-dice
model. It is the value of the function \(S\) on several elementary outcomes.

\begin{remarkbox}
A random variable often aggregates information. The sample space may distinguish
\((1,6)\) from \((6,1)\), but the sum random variable assigns both outcomes the
same value, namely \(7\).
\end{remarkbox}

\subsection*{Example: a payoff}

Random variables are not limited to counts. Suppose a game is played by rolling
one die. The player wins \(10\) units if the die shows \(6\), wins \(2\) units if
the die shows an even number different from \(6\), and loses \(3\) units
otherwise.

The sample space is
\[
\Omega=\{1,2,3,4,5,6\}.
\]
The payoff random variable \(X\) is defined by
\[
X(6)=10,
\]
\[
X(2)=X(4)=2,
\]
and
\[
X(1)=X(3)=X(5)=-3.
\]
Here the values of the random variable represent gains and losses. The
probability model describes the die; the random variable describes the numerical
consequence of each die result.

\subsection*{Indicator random variables}

One of the most useful random variables comes from a single event. Let \(A\) be
an event in \(\Omega\). The indicator of \(A\) is the function
\[
I_A(\omega)=
\begin{cases}
1, & \omega\in A,\\
0, & \omega\notin A.
\end{cases}
\]
This random variable records whether the event \(A\) occurred.

\begin{examplebox}
Roll one die and let
\[
A=\{2,4,6\}
\]
be the event that the result is even. The indicator \(I_A\) satisfies
\[
I_A(1)=0,\quad I_A(2)=1,\quad I_A(3)=0,\quad I_A(4)=1,\quad I_A(5)=0,\quad I_A(6)=1.
\]
It turns a yes-or-no event into a numerical variable.
\end{examplebox}

Indicators are simple, but they are powerful. Later, they will help us express
counts as sums. For example, the number of heads in \(n\) coin tosses can be
written as a sum of \(n\) indicators, one for each toss.

\subsection*{The inverse image of a numerical condition}

If \(X:\Omega\to\mathbb{R}\) is a random variable and \(B\subseteq\mathbb{R}\),
then the statement
\[
X\in B
\]
corresponds to an event in the original sample space:
\[
\{\omega\in\Omega:X(\omega)\in B\}.
\]
This set is called the inverse image of \(B\) under \(X\).

For example, if \(S(i,j)=i+j\) is the sum of two dice, then
\[
\{S=7\}=\{(1,6),(2,5),(3,4),(4,3),(5,2),(6,1)\}.
\]
The event \(\{S=7\}\) is not a new kind of object. It is a subset of the original
sample space.

\begin{remarkbox}
Every expression such as \(P(X=x)\), \(P(X\leq t)\), or \(P(X\in B)\) is really
a probability of an event in \(\Omega\). The random variable translates numerical
conditions into events.
\end{remarkbox}

\subsection*{From outcomes to values to distributions}

The workflow for random variables is:
\[
\text{experiment}\longrightarrow \Omega\longrightarrow X(\omega)\longrightarrow
\text{probabilities of values of }X.
\]
The first arrow constructs the sample space. The second defines the numerical
quantity of interest. The third asks how probability is transferred from
outcomes to numerical values.

This transferred probability law is called the distribution of the random
variable. The rest of the chapter develops the basic tools for describing such
laws in the discrete case.

\begin{theorembox}
A random variable does not create probability from nothing. It pushes the
probability already defined on \(\Omega\) onto numerical values. To understand a
random variable, first understand the sample space and then understand the
function defined on it.
\end{theorembox}
